\documentclass[../../main.tex]{subfiles}% 注意这里的文件路径不能用 ./main.tex ,否则用latexmk编译子文件会报错
\graphicspath{{\subfix{./image/}}{\subfix{../../image/}}} % 指定图片引用路径,后续可以直接使用图片文件名
% 如果图片存放在上级目录,则使用 ../../image/ 作为备用路径

% 例如:
% \begin{figure}[H]
% \centering
% \includegraphics[width=0.3\textwidth]{图.png}
% \caption{}
% \label{figure:图}
% \end{figure}
% 注意:上述\label{}一定要放在\caption{}之后,否则引用图片序号会只会显示??.

\begin{document}

\section{惯性定理}

\subsection{实二次型}

\begin{theorem}[实二次型的规范标准型]\label{theorem:实二次型的规范标准型}
设实二次型
\begin{align*}
f(x_1,x_2,\cdots,x_n)=d_1x_1^2 + d_2x_2^2 + \cdots + d_nx_n^2,
\end{align*}
这个实二次型对应的系数矩阵为$\boldsymbol{A}$ 是$n$ 阶实对称阵，则$\boldsymbol{A}$一定合同于下列对角阵:
\begin{align}
\mathrm{diag}\{1,\cdots,1;-1,\cdots,-1;0,\cdots,0\},\label{eq:8.3.3}
\end{align}
其中有 $p$ 个 $1$，$q$ 个 $-1$，$n - r$ 个零.进而,$f$一定可作变量替换得到
\begin{align}
y_1^2 + \cdots + y_p^2 - y_{p + 1}^2 - \cdots - y_r^2.\label{eq:8.3.2}
\end{align}
我们将\eqref{eq:8.3.2}式中的二次型称为$f$的\textbf{规范标准型}.
\end{theorem}
\begin{proof}
由\refthe{theorem:对称矩阵的合同标准型}，任意一个实对称阵 $\boldsymbol{A}$ 必合同于一个对角阵:
\begin{align*}
\boldsymbol{C}'\boldsymbol{A}\boldsymbol{C}=\mathrm{diag}\{d_1,d_2,\cdots,d_r,0,\cdots,0\},
\end{align*}
其中 $d_i\neq 0 (i = 1,\cdots,r)$. 注意到 $\boldsymbol{C}$ 是可逆阵，故 $r = \mathrm{r}(\boldsymbol{C}'\boldsymbol{A}\boldsymbol{C})=\mathrm{r}(\boldsymbol{A})$，即秩 $r$ 是矩阵合同关系下的一个不变量. 于是我们不妨设实对称阵已具有下列对角阵的形状:
\[
\boldsymbol{A}=\mathrm{diag}\{d_1,d_2,\cdots,d_r,0,\cdots,0\}.
\]

由\hyperref[lemma:初等合同变换]{初等合同变换}不难知道，任意调换 $\boldsymbol{A}$ 的主对角线上的元素得到的矩阵仍与 $\boldsymbol{A}$ 合同. 因此我们可把零放在一起，把正项与负项放在一起，即可设 $d_1>0,\cdots,d_p>0$; $d_{p + 1}<0,\cdots,d_r<0$. $\boldsymbol{A}$ 所代表的二次型为
\begin{align}
f(x_1,x_2,\cdots,x_n)=d_1x_1^2 + d_2x_2^2 + \cdots + d_rx_r^2.\label{eq:8.3.1}
\end{align}
再对上述二次型作变量替换,令
\[
\begin{cases}
y_1 = \sqrt{d_1}x_1,\cdots,y_p = \sqrt{d_p}x_p;\\
y_{p + 1} = \sqrt{-d_{p + 1}}x_{p + 1},\cdots,y_r = \sqrt{-d_r}x_r;\\
y_j = x_j (j = r + 1,\cdots,n),
\end{cases}
\]
则 \eqref{eq:8.3.1} 式变为
\begin{align*}
f = y_1^2 + \cdots + y_p^2 - y_{p + 1}^2 - \cdots - y_r^2.
\end{align*}
这一事实等价于说 $\boldsymbol{A}$ 合同于下列对角阵:
\begin{align*}
\mathrm{diag}\{1,\cdots,1;-1,\cdots,-1;0,\cdots,0\}.
\end{align*}
其中有 $p$ 个 $1$，$q$ 个 $-1$，$n - r$ 个零. 
\end{proof}

\begin{definition}
设 $f(x_1,x_2,\cdots,x_n)$ 是一个实二次型，若它能化为形如 \eqref{eq:8.3.2}式的形状，则称 $r$ 是该二次型的秩，$p$ 是它的\textbf{正惯性指数}，$q = r - p$ 是它的\textbf{负惯性指数}，$s = p - q$ 称为 $f$ 的\textbf{符号差}.
\end{definition}
\begin{remark}
显然，若已知秩 $r$ 与符号差 $s$，则 $p = \frac{1}{2}(r + s)$, $q = \frac{1}{2}(r - s)$. 事实上，在 $p,q,r,s$ 中只需知道其中两个数，其余两个数也就知道了. 由于实对称阵与实二次型之间的等价关系，我们将实二次型的秩、惯性指数及符号差也称为相应的实对称阵的秩、惯性指数及符号差.

因此\textbf{矩阵合同的充要条件是$p,q,r,s$中有两个数相同.}
\end{remark}

\begin{theorem}[惯性定理]\label{theorem:惯性定理}
证明\eqref{eq:8.3.2}式中的数$p$及$q=r-p$是两个合同不变量.即正惯性指数或负惯性指数相同的两个对称矩阵或二次型必合同.
这等价于:

设 $f(x_1,x_2,\cdots,x_n)$ 是一个 $n$ 元实二次型，且 $f$ 可化为两个标准型:
\begin{align*}
c_1y_1^2+\cdots + c_py_p^2 - c_{p + 1}y_{p + 1}^2 - \cdots - c_ry_r^2,\\
d_1z_1^2+\cdots + d_kz_k^2 - d_{k + 1}z_{k + 1}^2 - \cdots - d_rz_r^2,
\end{align*}
其中 $c_i>0$, $d_i>0$，则必有 $p = k$.
\end{theorem}
\begin{proof}
用反证法，设 $p>k$. 由\refthe{theorem:实二次型的规范标准型}不妨设 $c_i$ 及 $d_i$ 均为 $1$，因此
\begin{align}
y_1^2+\cdots + y_p^2 - y_{p + 1}^2 - \cdots - y_r^2 = z_1^2+\cdots + z_k^2 - z_{k + 1}^2 - \cdots - z_r^2.\label{eq:8.3.4}
\end{align}
又设
\[
\boldsymbol{x}=\boldsymbol{B}\boldsymbol{y},\ \boldsymbol{x}=\boldsymbol{C}\boldsymbol{z},
\]
其中
\[
\boldsymbol{x}=\begin{pmatrix}
x_1\\
x_2\\
\vdots\\
x_n
\end{pmatrix},\ \boldsymbol{y}=\begin{pmatrix}
y_1\\
y_2\\
\vdots\\
y_n
\end{pmatrix},\ \boldsymbol{z}=\begin{pmatrix}
z_1\\
z_2\\
\vdots\\
z_n
\end{pmatrix},
\]
于是 $\boldsymbol{z}=\boldsymbol{C}^{-1}\boldsymbol{B}\boldsymbol{y}$. 令
\[
\boldsymbol{C}^{-1}\boldsymbol{B}=\begin{pmatrix}
c_{11} & c_{12} & \cdots & c_{1n}\\
c_{21} & c_{22} & \cdots & c_{2n}\\
\vdots & \vdots & & \vdots\\
c_{n1} & c_{n2} & \cdots & c_{nn}
\end{pmatrix},
\]
则
\[
\begin{cases}
z_1 = c_{11}y_1 + c_{12}y_2 + \cdots + c_{1n}y_n,\\
z_2 = c_{21}y_1 + c_{22}y_2 + \cdots + c_{2n}y_n,\\
\cdots\cdots\cdots\\
z_n = c_{n1}y_1 + c_{n2}y_2 + \cdots + c_{nn}y_n.
\end{cases}
\]
因为 $p>k$，故齐次线性方程组
\[
\begin{cases}
z_1=c_{11}y_1 + c_{12}y_2 + \cdots + c_{1n}y_n = 0,\\
\cdots\cdots\cdots\\
z_k=c_{k1}y_1 + c_{k2}y_2 + \cdots + c_{kn}y_n = 0,\\
y_{p + 1} = 0,\\
\cdots\cdots\cdots\\
y_n = 0
\end{cases}
\]
必有非零解 ($n$ 个未知数，$n - (p - k)$ 个方程式). 令其中一个非零解为 $y_1 = a_1$, $\cdots$, $y_p = a_p$, $y_{p + 1} = 0$, $\cdots$, $y_n = 0$，把这组解代入 \eqref{eq:8.3.4} 式左边得到
\[
a_1^2+\cdots + a_p^2>0.
\]
但这时 $z_1 = \cdots = z_k = 0$，故 \eqref{eq:8.3.4} 式右边将小于等于零，引出了矛盾. 同理可证 $p<k$ 也不可能.
\end{proof}

\begin{theorem}\label{theorem:秩与符号差(或正负惯性指数)都是实对称阵在合同关系下的全系不变量}
秩与符号差 (或正负惯性指数) 是实对称阵在合同关系下的全系不变量.
\end{theorem}
\begin{proof}
由\hyperref[theorem:惯性定理]{惯性定理}知道，秩 $r$ 与符号差 $s$ 是实对称阵合同关系的不变量. 反之，若 $n$ 阶实对称阵 $\boldsymbol{A},\boldsymbol{B}$ 的秩都为 $r$，符号差都是 $s$，则它们都合同于
\[
\mathrm{diag}\{1,\cdots,1;-1,\cdots,-1;0,\cdots,0\},
\]
其中有 $p = \frac{1}{2}(r + s)$ 个 $1$，$q = \frac{1}{2}(r - s)$ 个 $-1$ 及 $n - r$ 个零，因此 $\boldsymbol{A}$ 与 $\boldsymbol{B}$ 合同. 对正负惯性指数的结论也同样成立. 
\end{proof}

\begin{definition}
令 $V$ 是实数域 $\mathbb{R}$ 上的 $n$ 维线性空间, $A$ 为 $n$ 阶实对称矩阵. 定义 $V$ 上的二次型为 $Q(x) = x^T A x$, 其中 $x \in V$ 是列向量.
设 $W$ 是 $V$ 的一个 $m$ 维子空间. \textbf{二次型 $Q$ 在 $W$ 上的限制}, 记为 $Q|_W$, 是指将定义域限制在 $W$ 上得到的二次型, 即 $Q|_W(x) = Q(x)$, 对任意 $x \in W$ 成立.

若在 $W$ 中选取一组基 $\{\bm{v}_1, \bm{v}_2, \cdots, \bm{v}_m\}$, 则 $Q|_W$的相伴实对称矩阵为 $G_W = (\bm{v}_i^T A \bm{v}_j)_{1 \leqslant i,j \leqslant m}$,即Gram矩阵.

并且子空间 $W$ 上的二次型 $Q|_W$ 的\textbf{正惯性指数}和\textbf{负惯性指数}, 分别定义为矩阵 $G_W$ 的正惯性指数 $p(G_W)$和负惯性指数 $q(G_W)$.
\end{definition}
\begin{remark}
由\hyperref[theorem:惯性定理]{惯性定理}知子空间 $W$ 上的二次型 $Q|_W$ 的正(负)惯性指数也与基$\{\bm{v}_1, \bm{v}_2, \cdots, \bm{v}_m\}$的选取无关.
\end{remark}

\begin{proposition}\label{proposition:子空间惯性指数一定小于全空间}
设$V$是$n$维线性空间,若全空间 $V$ 上二次型 $Q$ 的正(负)惯性指数为 $p$, 则对任意子空间 $W \subseteq V$, $Q|_W$ 的正(负)惯性指数 $p(W)$ 必定满足 $p(W) \leqslant p$.
\end{proposition}
\begin{proof}
仅证明正惯性指数的情形,负惯性指数同理可证.
已知 $V$ 上 $Q$ 的正惯性指数为 $p$, 则由\hyperref[theorem:惯性定理]{惯性定理}, 存在 $V$ 的一组标准正交基, 使得二次型在对应的标准正交坐标系下化为标准型
\begin{align*}
Q(x) = x_1^2 + \cdots + x_p^2 - x_{p+1}^2 - \cdots - x_n^2.
\end{align*}
于是 $V$ 可以分解为正定子空间 $V_+ = \mathrm{span}\{\bm{e}_1, \cdots, \bm{e}_p\}$ 与 $E = \mathrm{span}\{\bm{e}_{p+1}, \cdots, \bm{e}_n\}$ 的直和, 注意到对于任意 $x \in E$, 都有 $Q(x) \leqslant 0$, 并且 $\dim(E) = n-p$.

反证,假设存在某个子空间 $W$ 使得 $p(W) \geqslant p+1$. 那么在 $W$ 中必存在一个$p+1$维的子空间 $U$, 使得 $Q|_U$ 是严格正定的, 即对于任意非零向量 $x \in U$, 都有 $Q(x) > 0$.
因为 $U+E$ 必然是全空间 $V$ 的一个子空间, 所以 $\dim(U+E) \leqslant \dim(V) = n$.再根据\hyperref[theorem:交和空间维数公式]{交和空间维数公式}可得
\begin{align*}
\dim(U \cap E) = \dim(U) + \dim(E) - \dim(U+E) =(p+1) + (n-p)-\dim(U+E) \geqslant n+1 - n = 1.
\end{align*}
这个不等式意味着交集 $U \cap E$ 中必然存在一个非零向量 $x$. 然而, 因为 $x \in U$, 我们有 $Q(x) > 0$; 又因为 $x \in E$, 我们又得到 $Q(x) \leqslant 0$,显然矛盾!
\end{proof}

\begin{example}
对于 \(n\) 阶实对称矩阵 \(A\), 存在 \(y \in \mathbb{R}^n\) 使得 \(y^T A y > 0\). 证明: 对任意的 \(x \in \mathbb{R}^n\), 有
\begin{align}
(x^T A y)^2 \geqslant (x^T A x)(y^T A y). \label{eq:userin_1}
\end{align}
的充要条件是 \(x^T A x\) 的正惯性指数等于 \(1\).
\end{example}
\begin{remark}
对于充分性, 当 \(A\) 可逆时, 即 \(A\) 的正、负惯性指数分别为 \(1, n-1\) 时, \eqref{eq:userin_1} 式取等的充要条件是 \(\Delta = 0\), 对应 \(g\left(-\frac{x_1}{y_1}\right) = 0\), 由此可知 \(\frac{x_i}{y_i} = \frac{x_1}{y_1} \; (i=2,3,\cdots,n)\), 这意味着存在实数 \(k\), 使得 \(x = k y\).
\end{remark}
\begin{proof}
由于 \(A\) 为实对称矩阵, 所以存在可逆实矩阵 \(C\), 使得
\begin{align*}
C^T A C = \mathrm{diag}\{I_p, -I_q, O\}.
\end{align*}
注意到条件和结论在可逆线性替换$X=C^{-1}x,Y=C^{-1}y$下不变,故可不妨设$A=\mathrm{diag}\{I_p, -I_q, O\}$.
记 \(y = (y_1, y_2, \cdots, y_n)^T\), 对任意的 \(x \in \mathbb{R}^n\), 记 \(x = (x_1, x_2, \cdots, x_n)^T\), 则有
\begin{align*}
y^T A y = y_1^2 + \cdots + y_p^2 - y_{p+1}^2 - \cdots - y_{p+q}^2 > 0.
\end{align*}
由此可知 \(y_1, \cdots, y_p\) 不全为零,不妨设 \(y_1 \neq 0\).

{\heiti 必要性:}已知对任意的 \(x \in \mathbb{R}^n\), 有 \((x^T A y)^2 \geqslant (x^T A x)(y^T A y)\), 也就是对任意的 \(x_1, x_2, \cdots, x_n \in \mathbb{R}\), 有
\begin{align*}
(x_1 y_1 + \cdots + x_p y_p - x_{p+1} y_{p+1} - \cdots - x_{p+q} y_{p+q})^2 \geqslant (x_1^2 + \cdots + x_p^2 - x_{p+1}^2 - \cdots - x_{p+q}^2)(y_1^2 + \cdots + y_p^2 - y_{p+1}^2 - \cdots - y_{p+q}^2).
\end{align*}
若其中 \(p \geqslant 2\), 取 \(x_1 = y_2, x_2 = -y_1, x_3 = \cdots = x_n = 0\), 则此时上式等价于
\begin{align*}
0 \geqslant (y_1^2 + y_2^2)(y^TAy) > 0,
\end{align*}
矛盾！因此 \(p \leqslant 1\), 而由$y^TAy>0$知\(p \geqslant 1\), 因此 \(x^T A x\) 的正惯性指数为\(p = 1\).

{\heiti 充分性:}{\color{blue}证法一:}已知 \(x^T A x\) 的正惯性指数 \(p\) 等于 \(1\), 对任意的 \(x_1, x_2, \cdots, x_n \in \mathbb{R}\), 构造二次函数
\begin{align*}
g(t) = (t y_1 + x_1)^2 - \sum_{i=2}^{q+1} (t y_i + x_i)^2.
\end{align*}
其二次项系数 \(y_1^2 - y_2^2 - \cdots - y_{q+1}^2 = y^T A y > 0\), 因此这是开口向上的抛物线. 同时注意到
\begin{align*}
g\left(-\frac{x_1}{y_1}\right) = -\sum_{i=2}^{q+1} \left(-\frac{x_1}{y_1} \cdot y_i + x_i\right)^2 \leqslant 0.
\end{align*}
因此 \(g(t)\) 的判别式
\begin{align*}
\Delta = 4(x_1 y_1 - x_2 y_2 - \cdots - x_{q+1} y_{q+1})^2 - 4(x_1^2 - x_2^2 - \cdots - x_{q+1}^2)(y_1^2 - y_2^2 - \cdots - y_{q+1}^2) \geqslant 0.
\end{align*}
化简可得 \((x^T A y)^2 \geqslant (x^T A x)(y^T A y)\).

{\color{blue}证法二:}假设 \(x^T A x\) 的正惯性指数 \(p = 1\). 若存在向量 \(x\) 使得 \((x^T A y)^2 < (x^T A x)(y^T A y)\), 则关于 \(t\) 的二次函数
\begin{align*}
g(t) = (x + t y)^T A (x + t y)=(y^TAy)t^2+(x^TAy)^2t+x^TAx
\end{align*}
的判别式
$\Delta = 4[(x^TAy)^2-(x^TAx)(y^TAy)]<0$
又\(y^T A y > 0\),故$g(t)$恒大于0.
这意味着由 \(x, y\) 张成的二维子空间 \(W\) 是一个严格正定子空间, 其正惯性指数为 \(2\). 但根据\cref{proposition:子空间惯性指数一定小于全空间}, 子空间的正惯性指数不能超过全空间的正惯性指数 \(p\), 即有 \(2 \leqslant 1\), 矛盾!

{\color{blue}证法三:}
不妨设 \(x^TAx=x_1^2 - x_2^2 - \cdots - x_s^2 > 0\)，否则\eqref{eq:userin_1}式显然成立. 不妨设 \(x_1 y_1 \geqslant 0\)，否则用 \(-y\) 代替 \(y\)，现在由 Cauchy-Schwarz不等式得
\begin{align*}
x_1^2 y_1^2 &= \left[ \left( \sqrt{x_1^2 - x_2^2 - \cdots - x_s^2} \right)^2 + x_2^2 + \cdots + x_s^2 \right] \left[ \left( \sqrt{y_1^2 - y_2^2 - \cdots - y_s^2} \right)^2 + y_2^2 + \cdots + y_s^2 \right] \\
&\geqslant \left( \sqrt{x_1^2 - x_2^2 - \cdots - x_s^2} \sqrt{y_1^2 - y_2^2 - \cdots - y_s^2} + x_2 y_2 + \cdots + x_s y_s \right)^2.
\end{align*}
上式两边开方即得
\begin{align*}
x_1 y_1 - x_2 y_2 - \cdots - x_s y_s \geqslant \sqrt{x_1^2 - x_2^2 - \cdots - x_s^2} \sqrt{y_1^2 - y_2^2 - \cdots - y_s^2}.
\end{align*}
再对上式两边平方即得\eqref{eq:userin_1}式.
\end{proof}


\subsection{复二次型}

\begin{theorem}[复二次型的规范标准型]\label{theorem:复二次型的规范标准型}
设复二次型
\begin{align*}
f(x_1,x_2,\cdots,x_n)=d_1x_1^2 + d_2x_2^2 + \cdots + d_rx_r^2,
\end{align*}
这个复二次型对应的系数矩阵为$\boldsymbol{A}$ 是$n$ 阶复对称阵，则$\boldsymbol{A}$一定合同于下列对角阵:
\begin{align}
\mathrm{diag}\{1,\cdots,1;0,\cdots,0\},\label{eq:::8.3.3}
\end{align}
其中有 $r$ 个 $1$,并且$r=\mathrm{r}\left( A \right)$.进而,$f$一定可作变量替换得到
\begin{align}
z_1^2 + z_2^2 + \cdots + z_r^2.\label{eq:::8.3.2}
\end{align}
我们将\eqref{eq:::8.3.2}式中的二次型称为$f$的\textbf{规范标准型}.
\end{theorem}
\begin{proof}
因为复二次型
\begin{align*}
f(x_1,x_2,\cdots,x_n)=d_1x_1^2 + d_2x_2^2 + \cdots + d_rx_r^2
\end{align*}
必可化为
\begin{align*}
z_1^2 + z_2^2 + \cdots + z_r^2,
\end{align*}
其中 $z_i = \sqrt{d_i}x_i (i = 1,2,\cdots,r)$, $z_j = x_j (j = r + 1,\cdots,n)$. 所以结论得证.故复对称阵的合同关系只有一个全系不变量，那就是秩 $r$. 
\end{proof}














\end{document}